\documentclass[11pt]{article}
\usepackage[margin=1.1in]{geometry}
\usepackage{amsmath,amssymb,amsthm,mathtools,booktabs,array,tabularx,longtable}
\usepackage[dvipsnames]{xcolor}
\usepackage[colorlinks=true,linkcolor=MidnightBlue,citecolor=MidnightBlue,urlcolor=MidnightBlue]{hyperref}
\usepackage{fancyhdr}
\pagestyle{fancy}\fancyhf{}\renewcommand{\headrulewidth}{0pt}
\fancyfoot[R]{\thepage}
\fancyfoot[L]{\scriptsize\color{gray}Generated by the Physics Graph from approved parts of its store; not a deposited paper. Papers of record: \href{https://zenodo.org/communities/transport-papers}{Zenodo, Transport Papers}}
\fancypagestyle{plain}{\fancyhf{}\fancyfoot[R]{\thepage}\fancyfoot[L]{\scriptsize\color{gray}Generated by the Physics Graph from approved parts of its store; not a deposited paper. Papers of record: \href{https://zenodo.org/communities/transport-papers}{Zenodo, Transport Papers}}\renewcommand{\headrulewidth}{0pt}}
\newtheorem{theorem}{Theorem}[section]
\newtheorem{proposition}[theorem]{Proposition}
\newtheorem{lemma}[theorem]{Lemma}
\newtheorem{corollary}[theorem]{Corollary}
\newtheorem{conjecture}[theorem]{Conjecture}
\newtheorem{refutation}[theorem]{Refutation}
\newtheorem{prediction}[theorem]{Prediction}
\newtheorem{principle}[theorem]{Principle}
\newtheorem{observation}[theorem]{Observation}
\newtheorem{axiom}[theorem]{Axiom}
\theoremstyle{definition}
\newtheorem{definition}[theorem]{Definition}
\newtheorem{fact}[theorem]{Fact}
\newtheorem{assembled}[theorem]{Assembly}
\newtheorem{postulate}[theorem]{Postulate}
\newtheorem{premise}[theorem]{Premise}
\newtheorem{hypothesis}[theorem]{Hypothesis}
\newtheorem{remark}[theorem]{Remark}
\newtheorem{example}[theorem]{Example}
\theoremstyle{plain}
\setlength{\marginparwidth}{0.8in}
\newcommand{\changedmark}[2]{\marginpar{\raggedright\scriptsize\color{#1}#2}}
\providecommand{\Aa}{}\renewcommand{\Aa}{\mathcal{A}}
\providecommand{\Ad}{}\renewcommand{\Ad}{\widetilde{\mathrm{Ad}}}
\providecommand{\Aev}{}\renewcommand{\Aev}{\mathcal{A}^{+}}
\providecommand{\Aod}{}\renewcommand{\Aod}{\mathcal{A}^{-}}
\providecommand{\B}{}\renewcommand{\B}{\mathrm{B}}
\providecommand{\C}{}\renewcommand{\C}{\mathbb{C}}
\providecommand{\CAR}{}\renewcommand{\CAR}{\mathrm{CAR}}
\providecommand{\CS}{}\renewcommand{\CS}{{\mathbb{S}^{+}}}
\providecommand{\CSm}{}\renewcommand{\CSm}{{\mathbb{S}^{-}}}
\providecommand{\Cl}{}\renewcommand{\Cl}{\mathrm{Cl}}
\providecommand{\Con}{}\renewcommand{\Con}{\mathcal{C}}
\providecommand{\Dl}{}\renewcommand{\Dl}{\Delta}
\providecommand{\EN}{}\renewcommand{\EN}{E_{N}}
\providecommand{\End}{}\renewcommand{\End}{\mathrm{End}}
\providecommand{\F}{}\renewcommand{\F}{\mathbb{F}}
\providecommand{\G}{}\renewcommand{\G}{\Gamma}
\providecommand{\HH}{}\renewcommand{\HH}{\mathbb{H}}
\providecommand{\IS}{}\renewcommand{\IS}{I\sigma_3}
\providecommand{\J}{}\renewcommand{\J}{\mathbf{J}}
\providecommand{\LH}{}\renewcommand{\LH}{L_{H}}
\providecommand{\Mat}{}\renewcommand{\Mat}{\mathrm{M}}
\providecommand{\Mo}{}\renewcommand{\Mo}{M^{\circ}}
\providecommand{\Mp}{}\renewcommand{\Mp}{M^{+}}
\providecommand{\Msun}{}\renewcommand{\Msun}{M_{\odot}}
\providecommand{\Mt}{}\renewcommand{\Mt}{\widetilde{M}}
\providecommand{\Ncone}{}\renewcommand{\Ncone}{\mathcal{N}}
\providecommand{\Neff}{}\renewcommand{\Neff}{N_{\mathrm{eff}}}
\providecommand{\PCS}{}\renewcommand{\PCS}{\mathbb{P}{\mathbb{S}^{+}}}
\providecommand{\PT}{}\renewcommand{\PT}{\mathbb{PT}}
\providecommand{\Pin}{}\renewcommand{\Pin}{\mathrm{Pin}}
\providecommand{\Q}{}\renewcommand{\Q}{\mathbb{Q}}
\providecommand{\R}{}\renewcommand{\R}{\mathbb{R}}
\providecommand{\Ric}{}\renewcommand{\Ric}{\mathrm{Ric}}
\providecommand{\SBH}{}\renewcommand{\SBH}{S_{\mathrm{BH}}}
\providecommand{\Sc}{}\renewcommand{\Sc}{\mathbb{S}}
\providecommand{\Sch}{}\renewcommand{\Sch}{\mathrm{P}}
\providecommand{\Sig}{}\renewcommand{\Sig}{\Sigma}
\providecommand{\Som}{}\renewcommand{\Som}{S_\omega}
\providecommand{\Sp}{}\renewcommand{\Sp}{\mathbb{S}}
\providecommand{\Sph}{}\renewcommand{\Sph}{\mathrm{S}}
\providecommand{\Spi}{}\renewcommand{\Spi}{S_\pi}
\providecommand{\Spin}{}\renewcommand{\Spin}{\mathrm{Spin}}
\providecommand{\Sth}{}\renewcommand{\Sth}{S_{\mathrm{th}}}
\providecommand{\T}{}\renewcommand{\T}{\mathbb{T}}
\providecommand{\TH}{}\renewcommand{\TH}{T_{\mathrm{H}}}
\providecommand{\Tr}{}\renewcommand{\Tr}{\operatorname{Tr}}
\providecommand{\UHF}{}\renewcommand{\UHF}{\mathrm{UHF}}
\providecommand{\Z}{}\renewcommand{\Z}{\mathbb{Z}}
\providecommand{\Zc}{}\renewcommand{\Zc}{\mathrm{Z}}
\providecommand{\cmg}{}\renewcommand{\cmg}{\mathrm{cm^{2}\,g^{-1}}}
\providecommand{\dS}{}\renewcommand{\dS}{\mathrm{dS}}
\providecommand{\dsc}{}\renewcommand{\dsc}{\mathrm{dsc}}
\providecommand{\dtPG}{}\renewcommand{\dtPG}{dt}
\providecommand{\dtS}{}\renewcommand{\dtS}{dt_{\mathrm{S}}}
\providecommand{\ellv}{}\renewcommand{\ellv}{\ell}
\providecommand{\er}{}\renewcommand{\er}{e_r}
\providecommand{\et}{}\renewcommand{\et}{e_t}
\providecommand{\ev}{}\renewcommand{\ev}{\mathrm{ev}}
\providecommand{\gl}{}\renewcommand{\gl}{\mathfrak{gl}}
\providecommand{\hatV}{}\renewcommand{\hatV}{\widehat{V}}
\providecommand{\hb}{}\renewcommand{\hb}{\bar{h}}
\providecommand{\hf}{}\renewcommand{\hf}{\hat{f}}
\providecommand{\id}{}\renewcommand{\id}{\mathrm{id}}
\providecommand{\kpc}{}\renewcommand{\kpc}{\mathrm{kpc}}
\providecommand{\lP}{}\renewcommand{\lP}{\ell_{P}}
\providecommand{\mD}{}\renewcommand{\mD}{\mathbf{m}_{D}}
\providecommand{\mb}{}\renewcommand{\mb}{\bar m}
\providecommand{\mnu}{}\renewcommand{\mnu}{\Sigma m_{\nu}}
\providecommand{\nInf}{}\renewcommand{\nInf}{n_\infty}
\providecommand{\nO}{}\renewcommand{\nO}{n_0}
\providecommand{\ninf}{}\renewcommand{\ninf}{n_{\infty}}
\providecommand{\nz}{}\renewcommand{\nz}{n_{0}}
\providecommand{\proj}{}\renewcommand{\proj}{\operatorname{proj}}
\providecommand{\pv}{}\renewcommand{\pv}{\pi}
\providecommand{\rad}{}\renewcommand{\rad}{\mathrm{rad}}
\providecommand{\sgn}{}\renewcommand{\sgn}{\operatorname{sgn}}
\providecommand{\shA}{}\renewcommand{\shA}{\mathcal{A}}
\providecommand{\shC}{}\renewcommand{\shC}{\mathcal{C}}
\providecommand{\shE}{}\renewcommand{\shE}{\mathcal{E}}
\providecommand{\shJ}{}\renewcommand{\shJ}{\mathcal{J}}
\providecommand{\shK}{}\renewcommand{\shK}{\mathcal{K}}
\providecommand{\shL}{}\renewcommand{\shL}{L}
\providecommand{\shW}{}\renewcommand{\shW}{\mathcal{W}}
\providecommand{\sll}{}\renewcommand{\sll}{\mathfrak{sl}}
\providecommand{\so}{}\renewcommand{\so}{\mathfrak{so}}
\providecommand{\su}{}\renewcommand{\su}{\mathfrak{su}}
\providecommand{\supp}{}\renewcommand{\supp}{\operatorname{supp}}
\providecommand{\tr}{}\renewcommand{\tr}{\operatorname{tr}}
\providecommand{\uu}{}\renewcommand{\uu}{\mathfrak{u}}
\providecommand{\br}{}\renewcommand{\br}[1]{\langle #1\rangle}
\providecommand{\bra}{}\renewcommand{\bra}[1]{\langle #1|}
\providecommand{\file}{}\renewcommand{\file}[1]{\texttt{#1}}
\providecommand{\hp}{}\renewcommand{\hp}[2]{( #1 \,|\, #2 )}
\providecommand{\ip}{}\renewcommand{\ip}[2]{\langle #1, #2\rangle}
\providecommand{\ket}{}\renewcommand{\ket}[1]{|#1\rangle}
\providecommand{\rel}{}\renewcommand{\rel}[2]{D\bigl(#1 \,\|\, #2\bigr)}
\providecommand{\rev}{}\renewcommand{\rev}[1]{\widetilde{#1}}
\providecommand{\slashed}{}\renewcommand{\slashed}[1]{{#1}\mkern-10mu/}
\title{\bfseries Dark matter and the dark towers}
\author{Leon Butler\\ \normalsize\href{mailto:lbutler.math@gmail.com}{\texttt{lbutler.math@gmail.com}}\\ \normalsize ORCID: \href{https://orcid.org/0009-0006-4449-5617}{0009-0006-4449-5617}}
\date{Generated by the Physics Graph, 20 November 2026}
\begin{document}
\maketitle
\begin{abstract}
\normalsize
Astronomers find about 5.4 times as much dark matter as ordinary matter. In this framework the dark matter lives in separate towers that touch ours only through gravity.

This trail meets the towers, the mirror tower that copies ours, and the reason their matter stays dark and never decays. Under its premises the mirror accounts for one unit of the dark matter, a copy of ordinary matter to within half a percent.

A counting rule then matches the full observed ratio, but it is recorded as a fit, not a prediction. The walk ends by setting the dark sector, as the series builds it, against the data, where no arrangement yet passes every test.
\end{abstract}
\newpage
\noindent \textbf{Authorship and disclosure. }The founding idea of the series is the author's: that division by zero in a Clifford algebra should be read not as a prohibition but as a transport into a higher-dimensional algebra, and that the transported division need not be accompanied by an inverse operation. The direction of investigation and the critical review of the results are likewise the author's. The mathematical development and the drafting of the text were carried out in extended collaboration with Claude Fable 5 (Anthropic), an AI system.

\noindent\fbox{\parbox{\dimexpr\linewidth-2\fboxsep-2\fboxrule}{\textbf{This is not a deposited paper.} It was generated by the Physics Graph on 20 November 2026 and assembled from approved parts of its store. It follows results across several papers of the series; each is stated with its proof and credited to the paper that first proved it. The official versions of those papers are on Zenodo, in the Transport Papers community, \url{https://zenodo.org/communities/transport-papers}.}}
\newpage

\section*{Introduction}
The walk first sets out which dark towers exist and why they stay dark, then turns to how much matter they hold and what it is made of.

It closes with tests against observations.

\section{What we assume}

\begin{postulate}[E: equal mass per unit ball]\label{c:dark-matter.e-equal-mass-per-unit-ball}
The formula for the amount of dark matter needs to know how much matter each hidden sector holds. This assumption says it in a weaker form than before: there is one length at which a ball of that radius holds the same mass in every sector, each ball taken in that sector's own number of dimensions. It does not say that particles in different sectors weigh the same, or that there are equally many of them, which the earlier form of the rule assumed. It is still an assumption, not a result, and it is little more than the formula itself put as a single rule: together with the assumption about how thick a slice counts, it holds exactly when each sector has the share the formula gives it and the two slices are equally thick.
\end{postulate}

\begin{postulate}[Pop: the four odd towers are populated]\label{c:dark-matter.odd-towers-populated}
The formula for the amount of dark matter assumes that four of the hidden sectors actually contain matter. The series has no argument for that except the measurement itself: if they were empty, the dark matter would equal the ordinary matter, and it is about five times larger. So the assumption is read off from the very number the formula is meant to account for, and cannot count as separate evidence for it. This entry records it as an assumption in its own right.
\end{postulate}

\begin{postulate}[P1: the single frozen point]\label{c:dark-matter.p1-single-frozen-point}
This assumption says the origin of the universe happened once, in one way, for everything. Every hidden sector feels the same gravity, and every sector that can have a handedness received the same one. It is assumed, not derived. The assumption has three parts of different standing: that all sectors feel one gravity (partly provable, partly a choice), that they share one handedness (replaceable by another assumption, and not needed for the dark-matter total), and that one large number is common to all of them (assumed, with no derivation).
\end{postulate}

\begin{postulate}[P2: indifference, scoped to sector labels]\label{c:dark-matter.p2-indifference}
This assumption says the birth of the universe treats alike only things that are genuinely alike. The only hidden sectors that count as alike are pairs related by a reflection, so only those pairs are forced to match. Which member of the matching pair we live in is simply a matter of where we stand. What must match are the sector's fixed labels, such as its stored charge, not how hot it is: the assumption leaves the temperature of the mirror sector open, and it takes the number of modes the birth deposits to be finite.
\end{postulate}

\begin{postulate}[P3: nearest-dimension representatives]\label{c:dark-matter.p3-nearest-dimension-representatives}
Each hidden sector could in principle be realised with several different numbers of dimensions. This assumption picks, for each one, the number closest to our four, with a tie-break for the cases exactly in between. It is a choice the paper makes, not something it derives.
\end{postulate}

\begin{postulate}[R3: one cell per unshared direction]\label{c:dark-matter.r3-one-cell-per-unshared-direction}
When a hidden sector has a direction our world lacks, only a thin slice of it counts towards the gravity we feel, and this assumption fixes that slice at one unit of length. A mechanism is offered that gives this value, but it explains the choice rather than proving it. The unit is the same length for every sector: the radius of the ball in which another assumption (E) says every sector holds the same mass. It is not a sector's own particle spacing, which would bring the particles' masses and numbers back into the formula.
\end{postulate}

\begin{postulate}[SI: a single chiral spinor ideal at residues 0 and 4]\label{c:dark-matter.single-ideal-at-residues-zero-and-four}
The formula for the amount of dark matter counts five hidden sectors and leaves out two others. For that to be right, those two must hold no matter. The series arranges this by a choice in how a particle and its mass are defined in those two sectors, a choice made knowing that the formula needs it. With the other choices the formula gives about 6.2 or 8.3 in place of the measured 5.4. This entry records that choice as an assumption in its own right.
\end{postulate}

\begin{postulate}[Premise H1: a constituent's momentum]\label{c:dark-tower-matter.h1-momentum}
This is an assumption, not a result: a constituent's motion is read from a spinor in the tower's own algebra, linked to the constituent in the larger conformal algebra in a way that respects the tower's rotations. In our four-dimensional world this is already established; in three and five dimensions it is assumed. A later result shows that the formula itself is not a free choice: once the motion is read from such a spinor in a way that respects rotations and boosts, it is the only possibility up to an overall factor. What remains assumed is that the spinor exists, that the reading has that form, and that motions of several constituents add.
\end{postulate}

\begin{postulate}[Premise H2: gauge algebras as hierarchy automorphisms]\label{c:dark-tower-matter.h2-identification}
This is an assumption, not a result: the forces acting on a tower's matter are taken to be the symmetries of its smallest composite that can have mass. Our own tower's forces arise this way in the Standard-Model paper, at the level of matching observation, and the other towers are assumed to follow the same pattern.
\end{postulate}

\begin{postulate}[Premise U: the transposition rule (a definition)]\label{c:dark-tower-matter.u-transposition-rule}
This is an assumption, not a result: it fixes how the steps that give our own particle physics are repeated in the other towers. It says what a constituent is, which mixings of constituents and of their internal pairs count as symmetries, and that any square root of minus one must come from the tower's own algebra rather than being added from outside. It is best read as a definition: it says what the words mean in the other towers, and a different definition would give different results; nothing in nature could prove it wrong.
\end{postulate}

\begin{postulate}[Premise: the common ceiling coupling of the deposited dynamics]\label{c:synthesis.common-ceiling-coupling-premise}
The mirror world has exactly the same kinds of particles and forces as ours, but that does not by itself make its forces equally strong. For the mirror's matter to weigh exactly as much as ours, its forces must start out with the same strengths as ours at the highest energy the theory reaches. The series' indifference assumption, as now worded, only matches fixed labels such as a stored charge, not force strengths, so this assumption is added to say the two worlds start with equal strengths. It says nothing about how hot the mirror is, and nothing about the other hidden worlds. Its cost is plain: one of the 5.4 units of dark matter rests on it, and a small difference in starting strengths would change the mirror's share a great deal.
\end{postulate}

\section{Dark matter and the dark towers}

Each step below either says what the dark towers are or checks what they should do.

Real Clifford algebras repeat with period eight, and the series reads each repeating family as a tower. Along a tower the residue \[ p - q \pmod 8 \] is invariant, with $p$ and $q$ the numbers of positive and negative directions.

Taking each tower at the dimension nearest four (premise P3), exactly the residues $0$ and $4$ lack a single time direction, the two-time case excluded by importing the standard instability. That leaves six towers that admit matter: ours and five dark ones.

\begin{proposition}[Towers are residues, and the single-clock filter]\label{c:dark-matter.towers-are-residues-single-clock-filter}
Clifford algebras repeat with period eight, and the paper reads the eight repeating families as eight parallel towers, each labelled by a number from 0 to 7. If the dimension may vary freely, every label can host a world with exactly one direction of time, so a premise is needed to pick one representative per label. Taking the dimension nearest four (premise P3), two labels are left without a single time direction, one with none and one with two, the latter excluded by importing the standard instability of two-time physics; the other six can hold matter: our world and five dark ones. Which member of a mirror pair is called our world depends on a sign convention for time that premise P2 declares to carry no physics.
\emph{First proved in} T14, Prop 3.1 (\href{https://doi.org/10.5281/zenodo.22722953}{doi:10.5281/zenodo.22722953}).
\end{proposition}
\noindent\emph{Rests on:} \ref{c:dark-matter.imported-eight-towers-mod-8-clock}, \ref{c:dark-matter.p2-indifference}, \ref{c:dark-matter.p3-nearest-dimension-representatives}.

The closest dark tower is the mirror, reached by reversing the signature, so that time and space directions swap roles. Every input of the Standard-Model chain is invariant under that reversal, so the mirror carries its own copy of the derived Standard Model.

With a common ceiling and under the single frozen point (P1) and indifference (P2), \[ \Omega\_{$\mathrm{mirror}$} = \Omega\_{b} \] to half a percent, where $\Omega_{b}$ is the ordinary (baryonic) matter density. A cold mirror is, however, an exception to that same indifference, which would otherwise make the mirror as hot as ours.

\begin{theorem}[The mirror tower is a copy of this one]\label{c:dark-matter.mirror-tower-copy-one}
Swapping which directions count as time and which as space turns our world into a mirror world, and this result shows that every ingredient used to derive our particle physics survives the swap. So the mirror world has its own copy of the Standard Model as the series derives it (under Hypothesis (S)), with mirror photons and mirror neutrinos. If the two worlds share one cutoff scale and start with equal force strengths there (an assumption of its own, the common ceiling coupling), and under the framework's premises of a single frozen point (P1) and indifference (P2), which give an equal matter excess once that excess is read as a quantity fixed at the origin and not by temperature, it also has the same particle masses and the same inherited excess, which by a decision of 12 November 2026 is the excess of baryons minus leptons. Its matter then equals our ordinary matter to within half a percent, the gap coming from the two worlds' nuclear histories, provided baryons make up the same fraction of that excess there as here, which has not been shown. A later note (T14c) shows that the indifference principle, applied as widely as the paper first applied it, would also make the mirror world as hot as ours, which observations rule out; as now worded it leaves the mirror's temperature open, and a cold mirror then needs a birth out of equilibrium.
\emph{First proved in} T14, Thm 4.3 (\href{https://doi.org/10.5281/zenodo.22722953}{doi:10.5281/zenodo.22722953}).
\end{theorem}
\noindent\emph{Rests on:} \ref{c:dark-matter.imported-inputs-standard-model-chain}, \ref{c:dark-matter.type-residue-what-each-tower-can}, \ref{c:dark-matter.p1-single-frozen-point}, \ref{c:dark-matter.p2-indifference}, \ref{c:synthesis.common-ceiling-coupling-premise}.

A copy of our world raises the question of why we cannot see it.

Under P1, the sector decomposition and the wall-sequence theorem's hypotheses (a compact base and good couplings), the tower charge is exactly superselected, so dark-tower matter is exactly stable. That it only gravitates, at first order, comes from the series' acoustic no-go, conditional on the Indifference premise and premise (Q).

\begin{proposition}[Superselected, hence dark and exactly stable]\label{c:dark-matter.superselected-hence-dark-exactly-stable}
Assuming all the sectors share one spacetime metric set by a single frozen point (Premise P1), matter in the hidden sectors can affect us only through its weight, which is why it is dark; this part is carried over from earlier work, holds in the framework's spherically averaged description of the modes near horizons, and inherits two of that work's assumptions (the Indifference premise and premise (Q)). Given how the sectors are split, and provided the sheaf paper's conditions for its walls hold (a compact base and well-behaved couplings), this result argues that no process can change a particle's tower, so dark matter of this kind never decays; energy can still pass between the towers through gravity. The sectors met only once, at the universe's birth; a black hole can pass energy between them, but never a particle's sector label.
\emph{First proved in} T14, Prop 4.1 (\href{https://doi.org/10.5281/zenodo.22722953}{doi:10.5281/zenodo.22722953}).
\end{proposition}
\noindent\emph{Rests on:} \ref{c:dark-matter.imported-eight-towers-mod-8-clock}, \ref{c:dark-matter.imported-gravity-only-coupling}, \ref{c:dark-matter.residue-soldered-invariant-not-algebraic-one}, \ref{c:sheaf.infinite-wall-sequence}, \ref{c:dark-matter.p1-single-frozen-point}.

This gives the trail its first prediction.

Direct-detection searches should stay null permanently, because dark matter couples to ordinary matter through gravity alone at first order. The expectation is conditional on the single-frozen-point premise and on the acoustic no-go.

\begin{prediction}[Direct detection stays null]\label{c:pred.dark-matter-direct-detection-null}
Experiments that try to catch dark-matter particles bumping into ordinary atoms should never see anything. In this framework dark matter lives in separate towers that touch ours only through gravity, which is far too weak to register in such detectors. This expectation is conditional: it rests on the framework's premise of a single frozen point, and on a result carried over from earlier work that holds in a spherically averaged description and under two further assumptions of that work.
\emph{Kill condition:} Any confirmed dark-matter detection above the gravitational floor.
\end{prediction}

The next question is how much dark matter the towers hold. Under P1 to P4, with every tower's matter diluting as ours, the mirror supplies an asymmetric dark share of one $\Omega_{b}$ to half a percent.

That leaves a symmetric remainder of \[ 4.364\,\Omega\_{b} \] which is not derived.

\begin{proposition}[The population count: one share to half a percent, and a symmetric remainder]\label{c:dark-matter.population-count-one-share-half-percent}
Astronomers find about 5.4 times as much dark matter as ordinary matter. Assuming the three premises P1 to P3 (a single frozen point, indifference and nearest-dimension representatives; a fourth, a matter excess shared evenly across the oriented towers, is needed only for the mirror, where it follows from the mirror being a copy) and that the mirror world starts with the same force strengths as ours (the common ceiling coupling), and that every hidden sector's matter thins out with cosmic expansion as ours does, this result accounts for one of those units: a mirror copy of ordinary matter holding the same amount as ours to within half a percent. The other 4.4 units are not explained here; they would have to be a different, symmetric kind of matter in the other hidden sectors, felt by us only through gravity. The account is held at observation level on those premises.
\emph{First proved in} T14, Prop 4.5 (\href{https://doi.org/10.5281/zenodo.22722953}{doi:10.5281/zenodo.22722953}).
\end{proposition}
\noindent\emph{Rests on:} \ref{c:dark-matter.mirror-share-one-how-fragile}, \ref{c:dark-matter.mirror-tower-copy-one}, \ref{c:dark-matter.no-coupling-rule-odd-towers-follows}, \ref{c:dark-matter.projection-every-towers-content-dilutes-ones}, \ref{c:dark-matter.superselected-hence-dark-exactly-stable}, \ref{c:dark-matter.which-towers-can-carry-arrow}, \ref{c:dark-matter.p1-single-frozen-point}, \ref{c:dark-matter.p2-indifference}, \ref{c:dark-matter.p3-nearest-dimension-representatives}, \ref{c:synthesis.common-ceiling-coupling-premise}.

A separate counting rule attempts the whole ratio. Under premises R1 to R3, each populated tower holds one quantum of a common mass per unit ball of its own space, so a tower with $D$-dimensional space has share $V_{3}/V_{D}$, with $V_{D}$ the volume of the unit ball.

Summed over the towers, the dark-to-baryon ratio is \[ \frac{11\pi + 16}{3\pi} \approx 5.3643, \] against the observed $5.364 \pm 0.065$.

\begin{proposition}[The unit-ball rule]\label{c:dark-matter.unit-ball-rule}
This rule says that there is one length at which a ball of that radius holds the same mass in every tower, each ball taken in that tower's own space, and it compares the resulting density with ours. Towers with fewer or more spatial dimensions then contribute different amounts, fixed by the volumes of balls in two, three and four dimensions. With the assumptions that the four odd hidden sectors hold matter and that two further sectors hold none, the rule gives a dark-to-ordinary matter ratio of about 5.364, the observed value. It no longer assumes that particles in different sectors weigh the same or are equally numerous: only the mass per ball matters. The price of that weaker assumption is plain to state: together with the assumption about how thick a slice counts, it holds exactly when each sector has the share the formula gives it and the two slices are equally thick, so it is the formula put as one rule and not a separate reason for it. The relation was found by solving for the measured value, and many equally simple formulas land inside the measurement's error, so the agreement is a fit rather than a prediction.
\emph{First proved in} T14, Prop 7.6 (\href{https://doi.org/10.5281/zenodo.22722953}{doi:10.5281/zenodo.22722953}).
\end{proposition}
\noindent\emph{Rests on:} \ref{c:dark-matter.projection-every-towers-content-dilutes-ones}, \ref{c:dark-matter.towers-are-residues-single-clock-filter}, \ref{c:dark-matter.e-equal-mass-per-unit-ball}, \ref{c:dark-matter.odd-towers-populated}, \ref{c:dark-matter.r3-one-cell-per-unshared-direction}, \ref{c:dark-matter.single-ideal-at-residues-zero-and-four}.

A close match calls for caution rather than celebration.

The relation was found by fitting, and it is one of fourteen equally simple values within the error bar of one Planck column. It rests on Premise E (equal mass per unit ball) and R3, which together say no more than its shares, and on the costed clause that empties residues $0$ and $4$, so it is recorded as a fit at observation level, not a prediction.

\begin{conjecture}[The status of the rule]\label{c:dark-matter.status-rule}
The paper's formula for the amount of dark matter matches the measured value closely, and this assessment explains why that match should not be over-read. The formula was found by working backwards from the measurement, and many equally simple formulas fit just as well. It rests on two unproved premises (equal mass per unit ball, and the one-unit thickness; until 19 November 2026 it was stated with three, two of them about particles), which together say no more than the formula does, and on a further choice, made knowing the answer it needed, that two of the hidden families hold no matter; the match is to one Planck data set; another survey combination gives a lower value that no version of the formula can reach. So it is recorded as a fitted relation at observation level, not a prediction.
\emph{First proved in} T14, Obs 7.31 (\href{https://doi.org/10.5281/zenodo.22722953}{doi:10.5281/zenodo.22722953}).
\end{conjecture}

The counting rule says how much matter the dark towers hold, but not what that matter is made of. Building it by our own construction, a constituent's scalars depend on the tower: the real numbers in the three-dimensional towers, the complex numbers in ours and the mirror, and the quaternions in the five-dimensional towers.

\begin{theorem}[The constituent's field in each tower]\label{c:dark-tower-matter.constituents-field}
Each of the eight towers has its own kind of number for its basic constituents: real numbers, complex numbers or quaternions. This theorem reads off which from the classification of real Clifford algebras: complex for our tower and its mirror, real for the three-dimensional towers and one sterile tower, quaternionic for the five-dimensional towers and the other sterile tower. The answer is the same in the larger conformal algebra and in the tower itself.
\emph{First proved in} T14c, Thm 3.2 (\href{https://doi.org/10.5281/zenodo.24403592}{doi:10.5281/zenodo.24403592}).
\end{theorem}
\noindent\emph{Rests on:} \ref{c:dark-tower-matter.imported-clifford-structure}.

With those scalars, mass begins at two constituents in dimensions three and four and at one constituent in dimension five.

Under premise H2 and Premise U the gauge algebra is \[ \so(2) \oplus \so(2) \] in the three-dimensional towers and $\uu(3) = \su(3) \oplus \uu(1)$ in the five-dimensional towers. Reading a constituent's momentum from its spinor is established in our tower but a premise in three and five dimensions (premise H1).

\begin{theorem}[The gauge algebras of the odd towers]\label{c:dark-tower-matter.gauge-algebras}
Assuming that a tower's forces are the symmetries of its smallest composite that can have mass, as in the Standard-Model paper, and a rule for carrying that construction to other towers, this theorem finds the forces of the odd towers. The three-dimensional towers get two independent forces like electromagnetism. The five-dimensional towers get a colour force like our strong force together with one force like electromagnetism, nine force carriers in all.
\emph{First proved in} T14c, Thm 5.1 (\href{https://doi.org/10.5281/zenodo.24403592}{doi:10.5281/zenodo.24403592}).
\end{theorem}
\noindent\emph{Rests on:} \ref{c:dark-tower-matter.constituents-field}, \ref{c:dark-tower-matter.ladder}, \ref{c:dark-tower-matter.pairs}, \ref{c:dark-tower-matter.third-pairs-unit-five-dimensions}, \ref{c:dark-tower-matter.where-mass-begins}, \ref{c:dark-tower-matter.u-transposition-rule}, \ref{c:dark-tower-matter.h2-identification}, \ref{c:dark-tower-matter.h1-momentum}.

Those forces decide whether anything stands behind the counting rule's premise. In the three-dimensional towers the coupling has a dimension and nothing confines in the viable phase, and in the five-dimensional towers nothing confines at all, so the odd towers' masses are bare parameters.

Their number is not fixed either, so nothing in their content sets the masses or densities that premise E equates, for about $81$ percent of the dark matter. Their quanta are elementary, not the confined composites the dark-matter paper reads them as.

\begin{theorem}[Neither E nor the particle model has a mechanism in the odd towers]\label{c:dark-tower-matter.r1-has-no-mechanism-odd-towers}
The dark-matter paper's rule assumes every populated tower holds the same mass per unit ball, and its older particle picture assumes every dark particle has the same mass. In our tower and the mirror a particle's mass comes from a strong force, but this assembly-level argument shows that in the odd-dimensional towers, which hold about four fifths of the dark matter, masses are free parameters and nothing fixes how many particles there are. There, both assumptions are choices rather than consequences.
\emph{First proved in} T14c, Obs 9.4 (\href{https://doi.org/10.5281/zenodo.24403592}{doi:10.5281/zenodo.24403592}).
\end{theorem}
\noindent\emph{Rests on:} \ref{c:dark-matter.mirror-share-one-how-fragile}, \ref{c:dark-tower-matter.gauge-algebras}, \ref{c:dark-tower-matter.three-dimensional-content}, \ref{c:dark-tower-matter.anomalies}, \ref{c:dark-tower-matter.survival-sheets}, \ref{c:dark-tower-matter.five-dimensional-towers}, \ref{c:dark-tower-matter.imported-unit-ball-rule}, \ref{c:dark-tower-matter.population}, \ref{c:dark-matter.e-equal-mass-per-unit-ball}.

Abelian dark forces raise the question of dark radiation, which the early universe limits. The three-dimensional towers' radiation channels close, while the mirror and five-dimensional channels stay open.

The framework predicts \[ \Delta\Neff \approx 0 \] conditionally on those open channels closing, which holds if the descent's deposit is no hotter than estimated, and on a cold mirror, an exception to P2.

\begin{prediction}[No dark radiation]\label{c:pred.no-dark-radiation}
The afterglow of the early universe records how many kinds of light, fast-moving particles existed at the time. The framework predicts that the dark towers add none, provided the remaining ways of heating them stay closed, which holds if the descent left them no hotter than the paper estimates. It also needs the mirror world to have started cold, which a later note (T14c) shows is an unstated exception to the framework's own indifference principle; applied evenly, that principle would make the mirror as hot as our world and the prediction would fail. As the principle is now scoped, it applies to labels such as charges and not to temperatures, so it leaves the mirror's temperature free; a cold mirror then needs the descent to have left the towers out of equilibrium.
\emph{Kill condition:} One confirmed massless dark species.
\end{prediction}

The last step puts the whole dark sector, as the series builds it, against the data.

At observation level no configuration of the five dark towers as built passes every test. The least exposed needs four choices the series does not make, and it stays $1.24$ above Roux and Cline's limit.

\begin{conjecture}[The dark towers as built, against the data]\label{c:dark-tower-matter.towers-built}
At the level of observation, this table checks the five dark towers, as the series builds them, against radiation limits, colliding galaxy clusters, the mirror's astrophysics and the equal-mass rule. No arrangement passes every test as it stands. The best one needs four extra choices that the series does not make, and even then remains in mild tension with one published limit.
\emph{First proved in} T14c, Obs 11.1 (\href{https://doi.org/10.5281/zenodo.24403592}{doi:10.5281/zenodo.24403592}).
\end{conjecture}

\section{Results used from elsewhere}
The proofs above use the following results of the series, stated briefly; each is proved in the paper named.

\begin{fact}[Imported: the eight towers and the mod-8 clock]\label{c:dark-matter.imported-eight-towers-mod-8-clock}
The $(1,1)$ ascent preserves $p - q$ with $\Cl(p+k,q+k) \cong \Mat(2^{k}, \Cl(p,q))$, so the eight residues of $p - q$ modulo eight label eight parallel towers.
\emph{Source:} T14, Fact 2.5 (\href{https://doi.org/10.5281/zenodo.22722953}{doi:10.5281/zenodo.22722953}).
\end{fact}

\begin{fact}[Gravity-only coupling at first order, restated from the acoustic no-go]\label{c:dark-matter.imported-gravity-only-coupling}
Under the Indifference premise and premise (Q), matter that gravitates on this metric and shares no gauge charge with us acts on us through the metric-forming correlations only through its stress-energy, at first order, by the series' acoustic no-go restated; with P1 it does nothing else observable at that order.
\emph{Source:} T14, Fact 2.6 (\href{https://doi.org/10.5281/zenodo.22722953}{doi:10.5281/zenodo.22722953}).
\end{fact}

\begin{fact}[Imported: the inputs of the Standard-Model chain]\label{c:dark-matter.imported-inputs-standard-model-chain}
The Standard-Model chain of the series uses exactly a stated list of inputs, from the conformal completion and its conformal spinor ideal to the mass formula, pair-counting hypercharge, Hypothesis (S) and signature-free imports, a list carried as a hypothesis.
\emph{Source:} T14, Fact 4.2 (\href{https://doi.org/10.5281/zenodo.22722953}{doi:10.5281/zenodo.22722953}).
\end{fact}

\begin{fact}[The mirror share is one, and how fragile that is]\label{c:dark-matter.mirror-share-one-how-fragile}
Under the hypothesis of equal conversion (the mirror's ratio of baryon number to $B - L$ equals this tower's, which is not shown), the mirror share is exactly one because the mirror copies this tower (which discharges the democracy premise for it), but at one loop, over the ceiling range of Hypothesis (U), one extra Dirac flavour, or a ceiling coupling difference above about $0.1$, would move it far from one, and nothing follows for the odd towers.
\emph{Source:} T14, Prop 4.7 (\href{https://doi.org/10.5281/zenodo.22722953}{doi:10.5281/zenodo.22722953}).
\end{fact}

\begin{fact}[Three obstructions to a coupling rule for the odd towers]\label{c:dark-matter.no-coupling-rule-odd-towers-follows}
Three facts block a rule relating an odd tower's gauge coupling to ours: off dimension four the coupling is dimensionful with no fixed comparison length, no coupling value is derived anywhere in the series, and neither indifference nor the common ceiling coupling reaches towers off the reversal orbit; that any such rule is a further premise is a remark, not a theorem.
\emph{Source:} T14, Prop A.4 (\href{https://doi.org/10.5281/zenodo.22722953}{doi:10.5281/zenodo.22722953}).
\end{fact}

\begin{fact}[Projection: every tower's content dilutes as this one's]\label{c:dark-matter.projection-every-towers-content-dilutes-ones}
Under the sourcing rule read from P1, with any extra directions co-expanding or static and a missing direction filled by comoving sheets, every tower's projected matter dilutes as $a^{-3}$ and its radiation as $a^{-4}$, as in this tower.
\emph{Source:} T14, Lem 6.1 (\href{https://doi.org/10.5281/zenodo.22722953}{doi:10.5281/zenodo.22722953}).
\end{fact}

\begin{fact}[The radiation audit: what the type-by-residue theorem closes, and what stays open]\label{c:dark-matter.radiation-audit-what-theorem-closes-what}
The three-dimensional towers' radiation channels close (by Polyakov confinement, or by topological mass in T14c's phase (T)), while the mirror and five-dimensional channels stay open, so, with the rung scale of Hypothesis (U), $\Delta N_{\mathrm{eff}} \approx 0$ is a conditional stance: it needs those channels to close by a listed route and the deposit to be no hotter than the estimate.
\emph{Source:} T14, Prop 6.2 (\href{https://doi.org/10.5281/zenodo.22722953}{doi:10.5281/zenodo.22722953}).
\end{fact}

\begin{fact}[The residue is a soldered invariant, not an algebraic one]\label{c:dark-matter.residue-soldered-invariant-not-algebraic-one}
The residue is preserved by every isomorphism carrying vectors to vectors, but not by abstract algebra isomorphisms, under which residues $0$ and $2$, $3$ and $7$, and $4$ and $6$ coincide at equal dimension.
\emph{Source:} T14, Thm 3.5 (\href{https://doi.org/10.5281/zenodo.22722953}{doi:10.5281/zenodo.22722953}).
\end{fact}

\begin{fact}[Type by residue, and what each tower can carry intrinsically]\label{c:dark-matter.type-residue-what-each-tower-can}
A tower's spinor type (real, complex or quaternionic, with one or two modules) is a function of its residue alone, and it fixes the intrinsic symmetry group, $\Z_{2}$, $U(1)$ or $Sp(1)$, that each tower's ground algebra carries.
\emph{Source:} T14, Thm 3.4 (\href{https://doi.org/10.5281/zenodo.22722953}{doi:10.5281/zenodo.22722953}).
\end{fact}

\begin{fact}[Which towers can carry the arrow]\label{c:dark-matter.which-towers-can-carry-arrow}
Only the even-residue towers can be oriented, and of them, at their nearest-dimension representatives, only this one and the mirror are also causal, so they alone can carry the matter asymmetry; the oriented two-time fork would be a Kleinian copy of the Standard Model, kept closed by the imported two-time instability.
\emph{Source:} T14, Cor 3.11 (\href{https://doi.org/10.5281/zenodo.22722953}{doi:10.5281/zenodo.22722953}).
\end{fact}

\begin{fact}[Chern-Simons levels of the odd towers' content]\label{c:dark-tower-matter.anomalies}
The three-dimensional towers' bare Chern-Simons levels are forced into $\mathbb{Z} + \tfrac12$, breaking parity, while in the five-dimensional towers the abelian cubic level of the content under (S) is integral for every choice of mass signs.
\emph{Source:} T14c, Prop 6.5 (\href{https://doi.org/10.5281/zenodo.24403592}{doi:10.5281/zenodo.24403592}).
\end{fact}

\begin{fact}[The five-dimensional towers in cosmology and halos]\label{c:dark-tower-matter.five-dimensional-towers}
Under R3 the five-dimensional towers stay five-dimensional, confine nothing, meet the radiation bound for $T_D/T_\gamma < 0.085$ and scatter negligibly, and under U both may hold only the neutral species.
\emph{Source:} T14c, Comp 9.3 (\href{https://doi.org/10.5281/zenodo.24403592}{doi:10.5281/zenodo.24403592}).
\end{fact}

\begin{fact}[Imported: the structure of real Clifford algebras]\label{c:dark-tower-matter.imported-clifford-structure}
Imported: $\Cl^0(p,q) \cong \Cl(p,q-1)$, $\Cl(p,q)$ is a matrix algebra over $\R$, $\C$ or $\HH$ fixed by $p - q \bmod 8$, and a minimal left ideal's commutant is the division algebra $(fAf)^{\mathrm{op}}$.
\emph{Source:} C. Chevalley, The Algebraic Theory of Spinors, Columbia University Press (1954).; P. Lounesto, Clifford Algebras and Spinors, 2nd ed., Cambridge University Press (2001).; H. B. Lawson and M.-L. Michelsohn, Spin Geometry, Princeton University Press (1989)..
\end{fact}

\begin{fact}[Imported: the unit-ball shares]\label{c:dark-tower-matter.imported-unit-ball-rule}
Under E (equal mass per unit ball at one length) and R3 the mirror holds $18.6$ percent of the dark matter, each three-dimensional tower $24.9$ and each five-dimensional tower $15.8$, total share $(11\pi + 16)/(3\pi)$; the particle model R1 and R2 implies E and is not used.
\emph{Source:} T14c, Fact 2.4 (\href{https://doi.org/10.5281/zenodo.24403592}{doi:10.5281/zenodo.24403592}).
\end{fact}

\begin{fact}[The momentum-preserving mixings of constituents]\label{c:dark-tower-matter.ladder}
A mixing $U \in M_n(\F)$ of $n$ constituents preserves total momentum if and only if $UU^\dagger = 1$, giving $O(n)$, $U(n)$ or $Sp(n)$ in the three-, four- and five-dimensional towers.
\emph{Source:} T14c, Thm 3.7 (\href{https://doi.org/10.5281/zenodo.24403592}{doi:10.5281/zenodo.24403592}).
\end{fact}

\begin{fact}[The ladder pairs of the causal towers]\label{c:dark-tower-matter.pairs}
The three-dimensional towers have two real ladder pairs, while the four- and five-dimensional towers have a third pair that needs a complex unit, fixed up to sign where the field is $\C$.
\emph{Source:} T14c, Prop 4.1 (\href{https://doi.org/10.5281/zenodo.24403592}{doi:10.5281/zenodo.24403592}).
\end{fact}

\begin{fact}[How the odd towers could be populated]\label{c:dark-tower-matter.population}
At assembly level the odd towers' symmetric matter can fill E's density only by a protected deposited count or a weak-coupling freeze-out, and neither is derived.
\emph{Source:} T14c, Obs 9.5 (\href{https://doi.org/10.5281/zenodo.24403592}{doi:10.5281/zenodo.24403592}).
\end{fact}

\begin{fact}[Survival of sheet relics against annihilation]\label{c:dark-tower-matter.survival-sheets}
At the particle model's density, in decimal estimates, a symmetric sheet relic survives only for $P$ below $1$ to $2 \times 10^{-16}$, which needs $e^2/m$ above $7$ to $9 \times 10^7$ in phase (T), while phase (C) fails that density.
\emph{Source:} T14c, Comp 9.2 (\href{https://doi.org/10.5281/zenodo.24403592}{doi:10.5281/zenodo.24403592}).
\end{fact}

\begin{fact}[The third pair's unit in five dimensions]\label{c:dark-tower-matter.third-pairs-unit-five-dimensions}
At $r = 3$ the pseudoscalar $I_7$ is the unit that builds a third pair with $\uu(3)$ symmetry, other quaternionic units spoiling it, while at $r = 5$ every unit $u \in \HH$ does, the choices forming a sphere $S^2$.
\emph{Source:} T14c, Thm 4.2 (\href{https://doi.org/10.5281/zenodo.24403592}{doi:10.5281/zenodo.24403592}).
\end{fact}

\begin{fact}[The three-dimensional towers' content]\label{c:dark-tower-matter.three-dimensional-content}
The three-dimensional towers carry two species $A$, $B$, each charged under one $U(1)$, with no chirality, gauge-invariant Dirac masses and no baryon-like composite, and under (S) the constituent-level states fall off the charge lattice.
\emph{Source:} T14c, Prop 6.1 (\href{https://doi.org/10.5281/zenodo.24403592}{doi:10.5281/zenodo.24403592}).
\end{fact}

\begin{fact}[Where mass begins in each causal tower]\label{c:dark-tower-matter.where-mass-begins}
Mass begins at two constituents in dimensions three and four, where the reversion pairing is alternating, and at one constituent in dimension five, and total momenta are never spacelike.
\emph{Source:} T14c, Thm 3.5 (\href{https://doi.org/10.5281/zenodo.24403592}{doi:10.5281/zenodo.24403592}).
\end{fact}

\begin{fact}[The infinite wall sequence]\label{c:sheaf.infinite-wall-sequence}
For a compact wall and good couplings at every stage, the completion iterates indefinitely, giving an infinite sequence of tame shadow walls whose Clifford fibres ascend the Witt tower by one rung per wall.
\emph{Source:} T4, Cor 5.5 (\href{https://doi.org/10.5281/zenodo.22721669}{doi:10.5281/zenodo.22721669}).
\end{fact}

\begin{prediction}[No indirect dark-sector signals]\label{c:pred.no-indirect-dark-signals}
Conditionally on the single-frozen-point premise, and on the series' acoustic no-go (s-wave current algebra, Indifference premise, premise (Q)), indirect-detection surveys should see no dark-sector decay products, and no annihilation products above the gravitationally suppressed rate of a charge-asymmetric relic.
\emph{Kill condition:} Any confirmed indirect signal.
\end{prediction}

\section*{Conclusion}
Along the trail, the dark matter is the content of hidden towers felt only through gravity, with the mirror supplying one unit of it and the symmetric remainder left underived. As built, no arrangement of the dark towers yet passes every test against the data.

The trail rests on stated premises: a single frozen point (P1), indifference (P2), nearest-dimension representatives (P3), an evenly shared asymmetry (P4), Hypothesis (S) behind the mirror's copy, the counting premises E and R3, premises H1, H2 and U behind the odd towers' forces, and Hypothesis (U) behind the radiation audit. The predictions of null searches and no extra radiation are conditional on those premises, and a confirmed signal would count against them.

\end{document}